\chapter{Darstellung der möglichen Lösungswege}
\label{cha:Loesungswege}

Sie haben bereits die Problemstellung beschrieben und die fachlichen Grundlagen geschaffen. In diesem Kapitel werden mögliche Lösungswege beschrieben. Häufig gibt es verschiedene Lösungswege, daher gehen Sie in den folgenden wissenschaftlichen Schritten vor. Sie
\begin{enumerate}
	\item beschreiben die mögliche Lösungswege anhand eines Morphologischen Kastens, 
	\item entwickeln Bewertungskriterien für die Lösung,
	\item bewerten die sinnvollen Lösungsalternativen und
	\item wählen anhand der Auswahlkriterien einen Lösungsweg aus und begründen Ihre Wahl sachlich.
\end{enumerate}

%===================================================================================================
\section{Materialien}
Die, für einen Versuch notwendigen, Materialien werden hier präzise mit Artikelbezeichung und Anzahl in Form einer Auflistung aufgeführt. 
Folgende Angaben müssen unbedingt gemacht werden:

\begin{itemize}
	\item Prüfmaschine / Versuchseinrichtung mit Hersteller, Typbezeichnung, Inventarnummer usw.
	\item alle jeweils an der Maschine bzw. Einrichtung eingestellten Parameter
	\item alle Messgeräte, jeweils mit Hersteller, Typbezeichnung, Inventarnummer, eingestellten Parametern usw.
	\item verwendete Proben mit allen erforderlichen Angaben nach der jeweiligen DIN-, EN-, ISO-Norm bzw. andere Normen je nach Notwendigkeit
	\item bei nicht genormten Versuchen sinngemäße Angaben zur Probenform, zu Versuchsparametern, zu Temperaturen, physikalischen und chemischen Eigenschaften usw.
	\item sämtliche gemessenen Werte bzw. Versuchsergebnisse mit allen Parametern
	\item angewendete Auswertungsformeln mit vollständigen Quellenangaben usw. \cite[S.~24]{Hering2009}.
\end{itemize}

%================================================================================================
\section{Versuchsaufbau...}
"`...wird in diesem Kapitel beschrieben."'
Dieses Beispiel soll zeigen, dass der Dokumentteil-Titel sprachlich nicht in den darauffolgenden Text einbezogen werden darf.
Der erste Satz des nachfolgenden Textes muss deshalb ein vollständiger Satz sein, der jedoch die Sachaussage des Dokumentteil-Titels aufgreifen bzw. wiederholen darf.
Hinter der Dokumentteil-Nummer und dem Dokumentteil-Titel darf nie ein Satzzeichen wie Punkt, Doppelpunkt, Fragezeichen, Ausrufungszeichen usw. stehen \cite[S.~14~f.]{Hering2009}.
